% ===========================================================================
\section{Kraft--McMillan: when lengths fit}\label{sec:kraft-mcmillan}
\warmth{0}\quad\whisper{A decodable code spends a finite amount of room in the code tree.}

\begin{strand}
The Kraft sum turns a list of codeword lengths into a capacity test. Unique
decodability forces the sum to be at most one; conversely, every integer length
list that passes this test can be packed into a prefix code.
\end{strand}

\trigger{Use this test before trying to draw codewords: impossible length lists fail immediately.}

\block{The two directions}
\noindent\begin{tabularx}{\linewidth}{@{}l L@{}}
\toprule
{\color{indigo}direction} & \multicolumn{1}{l}{\color{madder}meaning}\\
\midrule
necessity & uniquely decodable $\Longrightarrow \sum_xD^{-\ell(x)}\le1$\\
sufficiency & integer lengths with $\sum_xD^{-\ell(x)}\le1$
  $\Longrightarrow$ a prefix code with exactly those lengths\\
\bottomrule
\end{tabularx}

\begin{theorem}[Kraft--McMillan inequality]\label{thm:kraft-mcmillan}
Let $C:\X\to\D^*$ be a $D$-ary uniquely decodable code, and write
$\ell_x=|C(x)|$. Then
\[
  \sum_{x\in\X}D^{-\ell_x}\le1.
\]
Conversely, if positive integers $(\ell_x)_{x\in\X}$ satisfy this inequality,
then there exists a $D$-ary prefix code $C$ such that $|C(x)|=\ell_x$ for every
$x\in\X$.
\end{theorem}

\block{Necessity: amplify by concatenating $n$ symbols}

\begin{proof}[Necessary direction]
Set
\[
  K=\sum_{x\in\X}D^{-\ell_x},\qquad
  \ell_{\min}=\min_x\ell_x,\qquad
  \ell_{\max}=\max_x\ell_x.
\]
For a source string $x^n=(x_1,\ldots,x_n)$, its encoded length is
$\ell(x^n)=\sum_{i=1}^n\ell_{x_i}$. Expanding the $n$th power gives
\[
  K^n
  =\sum_{x^n\in\X^n}D^{-\ell(x^n)}
  =\sum_{k=n\ell_{\min}}^{n\ell_{\max}}A_kD^{-k},
\]
where $A_k$ is the number of source strings in $\X^n$ whose encoded output has
length $k$.

Unique decodability makes these $A_k$ output strings distinct. There are only
$D^k$ strings of length $k$ over $\D$, so $A_k\le D^k$. Therefore
\[
  K^n
  \le \sum_{k=n\ell_{\min}}^{n\ell_{\max}}1
  =n(\ell_{\max}-\ell_{\min})+1.
\]
Taking $n$th roots and then letting $n\to\infty$ yields $K\le1$.
\end{proof}

\begin{yourturn}
Where exactly did unique decodability enter the proof?
\studentrecall{It makes the $A_k$ encoded strings distinct, hence $A_k\le D^k$.}
\ifshowanswers\else\workspace[2]\fi
\answerprose{It was used to turn source-string counting into output-string
counting: distinct source strings must give distinct encoded strings, so among
the $D^k$ possible strings of length $k$ we can have at most $A_k\le D^k$.}
\end{yourturn}

\block{Sufficiency: pack adjacent $D$-ary intervals}

\begin{proof}[Sufficient direction]
Relabel $\X=\{1,\ldots,M\}$ so that
$\ell_1\le\ell_2\le\cdots\le\ell_M$, and define
\[
  r_m=\sum_{i=1}^{m-1}D^{-\ell_i},\qquad 1\le m\le M.
\]
The Kraft inequality gives $0\le r_m<1$. Since $\ell_i\le\ell_m$ for $i<m$,
$D^{\ell_m}r_m$ is an integer, so $r_m$ has a terminating base-$D$ expansion
with at most $\ell_m$ digits. Let $C(m)$ be its first $\ell_m$ digits, padding
with zeros when necessary.

To see that the resulting code is prefix-free, take $m<n$. By construction,
\[
  r_n-r_m=\sum_{i=m}^{n-1}D^{-\ell_i}\ge D^{-\ell_m}.
\]
But every base-$D$ fraction whose first $\ell_m$ digits equal $C(m)$ lies in
the half-open interval $[r_m,r_m+D^{-\ell_m})$. Thus $r_n$ cannot begin with
$C(m)$, so $C(m)$ is not a prefix of $C(n)$. Hence $C$ is a prefix code with
the prescribed lengths.
\end{proof}

\begin{example}[a tight binary packing]
For lengths $(1,2,3,3)$,
\[
  2^{-1}+2^{-2}+2^{-3}+2^{-3}=1.
\]
The construction uses left endpoints $0,\tfrac12,\tfrac34,\tfrac78$, whose
binary expansions give the prefix code $0,10,110,111$.
\end{example}

\begin{yourturn}
Could a binary uniquely decodable code have lengths $(1,2,2)$? Compute the
Kraft sum and state the verdict.
\[
  2^{-1}+2^{-2}+2^{-2}=\answerin[1.2cm]{1}.
\]
\studentrecall{The sum is $1$, so the lengths are feasible; for example $0,10,11$.}
\answerprose{Yes. The Kraft sum is $1$, so the converse constructs a prefix
code; one choice is $0,10,11$.}
\end{yourturn}

\vspace{1.2cm}
\recall{Why does the $n$th-root limit force the Kraft sum to be at most one?}
